% !TeX root = ../main.tex
\section{Misleading Setup}\label{sec:pitfall}
\rubric{AI-use transparency and pitfall analysis}{15}{Honest AI-use note and a meaningful misleading-setup
analysis (with \cref{sec:ai-use}).}

\begin{hypothesis}
  Dropping \code{O_DIRECT} lets the OS page cache serve repeated reads of the same region from RAM after
  the first pass, so a buffered benchmark will report throughput and latency that describe DRAM, not the
  SSD---orders of magnitude ``better'' than the real device, and entirely fictitious as a storage result.
\end{hypothesis}

\begin{knobs}
  Pitfall     & Buffered I/O measuring the page cache instead of the device \\
  Wrong setup & \qty{4}{\kibi\byte} sequential reads over a \qty{64}{\mebi\byte} region, \emph{no} \code{O_DIRECT} \\
  Right setup & Identical workload \emph{with} \code{O_DIRECT} (page cache bypassed) \\
  Script      & \code{SSD_Performance_Profiling/scripts/run.sh pitfall} \\
\end{knobs}

\begin{table}[tbp]
  \centering\small
  \caption{Same \qty{4}{\kibi\byte} sequential-read workload, direct vs.\ buffered (medians). The only
  difference is the \code{O_DIRECT} flag.}\label{tab:pitfall}
  \begin{tabularx}{\linewidth}{@{}Xrrrr@{}}
    \toprule
    Setup & IOPS & MiB/s & Mean (\si{\micro\second}) & p99 (\si{\micro\second}) \\
    \midrule
    \code{O_DIRECT} (correct) & \num{86201}   & \num{337}  & \num{11.57} & \num{21.3} \\
    buffered (page cache)     & \num{2057063} & \num{8035} & \num{0.45}  & \num{0.91} \\
    \bottomrule
  \end{tabularx}
\end{table}

\subsection{Explanation}
The buffered run reports \qty{2.06}{} million IOPS at \qty{8.0}{\gibi\byte\per\second} and sub-microsecond
latency---a $24\times$ inflation over the correct \qty{86}{k} IOPS / \qty{337}{\mebi\byte\per\second}
figure. None of it is a storage measurement: after the first pass the \qty{64}{\mebi\byte} region lives
entirely in the page cache, so every subsequent read is a memcpy from RAM. \qty{8}{\gibi\byte\per\second}
is DRAM bandwidth, and \qty{0.45}{\micro\second} is a cache-hit memory latency, not an NVMe access. This is
the single easiest way to publish a fraudulently fast SSD number, and an AI-generated benchmark falls into
it readily because omitting \code{O_DIRECT} is the \emph{default} behaviour of \code{open()}. Three cheap
detectors flag it: (1) the throughput exceeds the device's interface ceiling and even its own
QD=1/QD=128 direct numbers by more than an order of magnitude; (2) sub-microsecond ``storage'' latency is
physically impossible for NAND flash; and (3) the number scales with the \emph{region size} (shrinks once
the region exceeds RAM), which a true device result must not. Our methodology avoids it structurally:
\code{O_DIRECT} is the default for every real run (the buffered path exists only to produce this
comparison), and the validation step (\cref{sec:validation}) asserts the direct-vs-buffered gap to prove
the cache is bypassed. The same trap underlies ``wrong effective queue depth'': if the async engine never
keeps iodepth requests in flight, a QD=32 label can really be a QD=2 result---which is why we verify the
effective depth with the Little's Law table (\cref{tab:littles-law}) rather than trusting the requested
value.
